\begin{figure}
    \begin{tikzpicture}[node distance=20pt, every node/.style={fill=white, font=\sffamily, scale=0.6}, align=left]
        \node (start) [code] {
            \begin{lstlisting}
class
    EXAMPLE
feature
    min (x, y: INTEGER): INTEGER
        -- Header comment
        do
            Result := x
        ensure
            res: (Result = x or Result = y)
                and Result <= x and Result <= y
        end

    max (x, y: INTEGER): INTEGER
        -- Header comment
        do
            Result := x
        ensure
            res: (Result = x or Result = y)
                and Result >= x and Result >= y
        end
end
\end{lstlisting}
        };
        \node (start-description) [description, right of=start, xshift=135pt] {Initial feature in class};
        \draw[->] (start-description) -- (start);

        \node (prompt) [code, below of=start, yshift=-190pt] {
\begin{lstlisting}
class
    EXAMPLE
feature
    -- <task>
    min (x, y: INTEGER): INTEGER
        -- Header comment
        do
            Result := x
        ensure
            res: (Result = x or Result = y)
                and Result <= x and Result <= y
        end
        -- AutoProof counterexample

    -- <task>
    max (x, y: INTEGER): INTEGER
        -- Header comment
        do
            Result := x
        ensure
            res: (Result = x or Result = y)
                and Result >= x and Result >= y
        end
        -- AutoProof counterexample
end
\end{lstlisting}
        };
        \draw[->, thick] (start) -- (prompt);
        \node (prompt-description) [description, right of=prompt, xshift=130pt] {Generated prompt};
        \draw[->] (prompt-description) -- (prompt);

        \node (llm-candidate) [code, below of=prompt, yshift=-240pt] {
            \begin{lstlisting}
class
    EXAMPLE
feature
    -- <task>
    min (x, y: INTEGER): INTEGER
        -- Header comment
        do
            if x < y then
                Result := x
            else
                Result := y
            end
        ensure
            res: (Result = x or Result = y)
                and Result <= x and Result <= y
        end
        -- AutoProof counterexample

    -- <task>
    max (x, y: INTEGER): INTEGER
        -- Header comment
        do
            if x < y then
                Result := y
            else
                Result := x
            end
        ensure
            res: (Result = x or Result = y)
                and Result >= x and Result >= y
        end
        -- AutoProof counterexample
end
\end{lstlisting}
        };
        \draw[->, thick] (prompt) -- (llm-candidate);
        \node (llm-candidate-description) [description, right of=llm-candidate, xshift=140pt] {Generated llm-candidate};
        \draw[->] (llm-candidate-description) -- (llm-candidate);

        \node (verification-result) [code, below of=llm-candidate, yshift=-150pt] { Verification output };
        \draw[->, thick] (llm-candidate) -- (verification-result);
        \draw[->, thick] (verification-result.west) -| ([xshift=-30pt] prompt.west)
            node[near start, below] {Reset failing features}
            node[near start, above] {Some feature fails} -- (prompt.west);
        \node (initial-feature-after-edits) [code, below of=verification-result, yshift=-140pt] {
            \begin{lstlisting}
class
    EXAMPLE
feature
    min (x, y: INTEGER): INTEGER
        -- Header comment
        do
            if x < y then
                Result := x
            else
                Result := y
            end
        ensure
            res: (Result = x or Result = y)
                and Result <= x and Result <= y
        end

    max (x, y: INTEGER): INTEGER
        -- Header comment
        do
            if x < y then
                Result := y
            else
                Result := x
            end
        ensure
            res: (Result = x or Result = y)
                and Result >= x and Result >= y
        end
end
\end{lstlisting}
        };
        \draw[->, thick] (verification-result) -- (initial-feature-after-edits)
            node[midway, right] {All features verify};
        \node (initial-feature-after-edits-description) [description, right of=initial-feature-after-edits, xshift=130pt] {Feature after edits};
        \draw[->] (initial-feature-after-edits-description) -- (initial-feature-after-edits);
    \end{tikzpicture}
    \caption{Dataflow of the class by class correction command}\label{dataflow-class-by-class-correction}
\end{figure}